\begin{apient}
Breakpoint::ptr Breakpoint::const_ptr
\end{apient}

\apidesc{
The Breakpoint::ptr and Breakpoint::const\_ptr types are respectively a pointer
and a const pointer to a Breakpoint object. These pointers are shared
pointers, and the underlying Breakpoint object will be automatically clean when
there are no more references to it. ProcControlAPI will automatically
maintain at least one reference to any Breakpoint that is installed in a target
process.
}

\begin{apient}
static const int BP_X = 1 static const int BP_W = 2 static const int BP_R = 4
\end{apient}

\apidesc{
These constant values are used to set execute, write and read bits on hardware
breakpoints.
}

\begin{apient}
Breakpoint::ptr newBreakpoint()
\end{apient}

\apidesc{
This function creates a new software Breakpoint object and returns it. The
Breakpoint is not inserted into a Process until it is passed to
Process::addBreakpoint().
}

\begin{apient}
Breakpoint::ptr newTransferBreakpoint(Dyninst::Address ctrl_to)
\end{apient}

\apidesc{
This function creates a new control transfer software breakpoint. Upon
resumption after executing this Breakpoint, control will resume at the address
specified by the ctrl\_to parameter.
}

\begin{apient}
Breakpoint::ptr newHardwareBreakpoint(unsigned int mode, unsigned int size)
\end{apient}

\apidesc{
This function creates a new hardware breakpoint. The mode parameter is a
bitfield that contains an OR combination the values BP\_X, BP\_W and BP\_R.
These control whether the breakpoint will trigger when its target address is
executed, written or read.
The size parameter specifies a range of memory that this breakpoint monitors.
If memory is accessed between the target address and target address + size,
then the breakpoint will trigger.
}

\begin{apient}
bool isCtrlTransfer() const
\end{apient}

\apidesc{
This function returns true if the Breakpoint is a control transfer breakpoint,
and false if it is a regular Breakpoint. 
}

\begin{apient}
Dyninst::Address getToAddress() const
\end{apient}

\apidesc{
If this Breakpoint is a control transfer breakpoint, then this function returns
the address to which it transfers control. If this Breakpoint is not a
control transfer breakpoint, then this function returns 0.
}

\begin{apient}
void setData(void *data) const
\end{apient}

\apidesc{
This function sets the value of an opaque handle that is associated with each
Breakpoint. The opaque handle can be any value, and it can be retrieved with
the getData function. 
}

\begin{apient}
void *getData() const
\end{apient}

\apidesc{
This function returns the value of the opaque handled that is associated with
this Breakpoint. 
}

\begin{apient}
void setSuppressCallbacks(bool val)
\end{apient}

\apidesc{
This function can be used to suppress callbacks stemming from a specific
breakpoint when called with val set to true value. All other effects from
this breakpoint will still occur, but it will not generate a callback. By
default callbacks occur from every breakpoint.
}

\begin{apient}
bool suppressCallbacks() const
\end{apient}

\apidesc{
This function returns true if callbacks have been suppressed for this breakpoint
from Breakpoint::setSuppressCallbacks and false otherwise.
}
